\chapter{Against Ourselves}
The precise moment at which Rebellion turns us against ourselves is about a fourth of the way through the film. Up until that point, it has depicted the happy world that, as any viewer with a trace of empathy must concede, the characters have more than earned. Throughout the series, we followed these young women, suffered with them, hoped against hope that they would be able to find some form of happiness.

In the end of the series, arguably, they did, but there is little denying that the ending to the series is bittersweet. Homura is alone, the only one who remembers Madoka. Madoka is gone forever, never born to begin with. The rest of the magical girls still fight, still suffer, still sink into the uttermost depths of despair to become witches—but are mercy-killed by Madoka just as they do. 

Sayaka still died for a boy who barely noticed she was there. Mami and Kyoko are active as magical girls, so we can presum Mami's parents are still dead and Kyoko's family still perished in a murder-suicide. 

The end of the series was an honest ending, not a happy one. It depicted the creation of a new, better world, but far from a flawless one. 

That flawless world is what we see in the first segment of Rebellion. All five magical girls are alive and working as a team. Their interpersonal difficulties are reduced to flirtatious teasing between Sayaka and Kyoko. The psychic damage of Homura's time-travel shenanigans seems healed: Homura is back to her shyer, less confident, but more pleasant and cheerful glasses-wearing pigtailed self, and Madoka is both more cheerful and more confident, more like the version Homura first met at the beginning of Episode 10. 

The Nightmares are almost laughable as a threat. If Hitomi's Nightmare is anything to go by, they pose no physical threat to the girls, don't torture them psychologically, and can be reduced to literal moe-blobs. What's more, they release a massive abundance of Soul Gem-cleansing light when killed, which as I've noted before not only permits, but encourages, the girls to work together, and in addition provides more than enough energy to keep them from blackening their Soul Gems and dying. Instead, the girls get to be magically powerful and visually impressive, fighting as a team against just enough difficulty to feel useful without ever experiencing the horrors of the series. 

This is what we, collectively, as an audience, wanted. Oh, most of us understood that the ending as it stood was probably aesthetically better, but enough fanfiction by those too inexperienced to know better or too invested to care exists to make it clear: we wanted better for these girls. And here the movie comes, and gives us exactly what we asked for—until Homura starts to figure it out. 

Like Paradise Lost before it (as discussed in Chapter 19), the show tricks us into rooting for someone who is trying to destroy our paradise. Homura knows this happy world is untrue, and therefore we know that by investigating it she will destroy it. From her initial conversation with Kyoko, the world becomes less and less realistic, until by the time the two realize they are trapped in the city the world is an abstraction of red field and white lines, the bus the only recognizable object. Soon after, Homura becomes the familiar glasses-less, straight-haired, darkly stoic girl we remember from the series, and the familiar sight excites us even as it means the happy world is deteriorating still faster. 

Later, we see the battle teased throughout the first three episodes of the series, as Mami and Homura come to blows. The resulting battle is visually stunning, as Homura and Mami both employ their respective powers and extensive arsenals to the fullest. It is exciting, dramatic, well-animated and -scored—and horribly, horribly wrong. As a set piece, it is a long sequence that advances the plot little, the characters and themes not at all; it is exciting, but blatantly gratuitous, a pure piece of audience pandering of the sort the show deliberately shied away from most of the time. And then Homura shoots herself in the head, and Mami dissolves into ribbons, the pandering turned suddenly to horror. 

Getting what we want is a disappointment and leads to horror. Nowhere is this as clear as in the film's climax, when Homura and Madoka are reunited and it all goes horribly wrong, resulting in a world where all the girls are free and alive and Madoka doesn't have to be a magical girl—a corrupt world ruled by a demonic demiurgic Homura who is holding Madoka prisoner. 

We bought our tickets, sealed our contracts, and got our wishes, and they turned to ashes around us. Desire leads inevitably to suffering. 

Why? Because we might wish for happiness, but we need truth. This is not to say that despair is truer than happiness, but rather that the truth of Madoka Magica is entropy and the inevitability of decay, and the series has consistently equated physical entropy and decay to the feelings of depression and despair. To end straightforwardly, uncomplicatedly happily, to give us what we wish for without corrupting it or snatching it away is to deny itself. 

So the film forces us to reject our own desires for the series. Those who revel in its darkness and spiky difficulty must endure being pandered to with fanservice, pushing them to deny their own fandom. Those who embrace the fanservice must face where it leads. Both must deal with the deeply ambiguous final arc of the film, as Homura creates a world simultaneously darker and brighter than the world of the series, yet more coherent than the dream-world of the film.

Thus, the series places the viewer into the position of the magical girls. Pursuing our desires for the series leads to it becoming tragic. Our wishes transform into curses as down the spiral we go, until we find ourselves, at the climax of the film, wishing for Homura the witch to tear apart the world—and then when she does, we must live with the reality created by that wish.

By turning us against ourselves, and showing how our wishes for the series betray us, the film makes one last effort to push empathy onto us. Like the series in its first few episodes, it offers spectacle and fanservice to draw us in, and then, once the trap is baited, it makes us feel for the characters. Even more so, however, it makes us feel as the characters—empathy as opposed to sympathy—by placing us into a situation analogous to theirs. That moment of confusion, of alienation, of wrongness when Homura pulls Madoka apart? That is a small taste of what it feels like to be a magical girl.

I said above that this is a series about entropy and decay, depression and despair, and it is. But it's easy to forget that it's about other things, too, and by turning us against ourselves it reminds us of those other things.

This isn't just about Buddhism, or German literature, or the magical girl genre. It isn't just about entropy and suffering, or just about thematic complexity or the possible psychological issues of its implied, gestalt author. It isn't even just about characters, blobs of light and color created by animators and voiced by actors. It's also about us, our fears and hopes, our wishes and despair, our capacity for empathy.

In the end, as in the beginning, Puella Magi Madoka Magica is a story about people.
